% !TEX root = ../main.tex
\chapter{Loss Functions}
\label{ch:losses}

The loss function is the scalar objective that gradient descent
minimises. The choice of loss encodes the modelling assumption:
mean-squared error implies a Gaussian likelihood; cross-entropy
implies a categorical likelihood; Huber implies robustness to
outliers. \Lucid\ provides the standard loss family at
\code{lucid.nn} and \code{lucid.nn.functional}, mirroring
\refframework's surface.

This chapter catalogues the loss functions with their
mathematical definitions, the numerical-stability tricks for
each, the reduction conventions, and the gradient flow into the
model.

\section{The Reduction Convention}
\label{sec:losses:reduction}

Every \Lucid\ loss function accepts a \code{reduction} argument:

\begin{itemize}
  \item \code{'mean'} (default): average over all elements.
  \item \code{'sum'}: sum over all elements.
  \item \code{'none'}: no reduction; return per-element loss.
\end{itemize}

The \code{'mean'} default matches \refframework. For some losses
(notably \code{cross\_entropy}), the convention is to mean over
the batch, sum over the per-sample loss dimensions.

\subsection{Why the reduction matters}
\label{ssec:losses:why_reduction}

The gradient with respect to the parameters scales with the
reduction:
\[
\nabla_\theta \mathcal{L}_{\text{mean}} \;=\; \frac{1}{N} \nabla_\theta \mathcal{L}_{\text{sum}}.
\]
Switching from \code{'mean'} to \code{'sum'} at constant learning
rate scales the effective gradient by $N$, equivalent to
multiplying the learning rate by $N$. This is the source of the
``my training broke when I changed reduction'' bug class.

\section{Regression Losses}
\label{sec:losses:regression}

\subsection{Mean Squared Error (L2)}
\label{ssec:losses:mse}

\[
\mathcal{L}_{\text{MSE}}(y, \hat y) \;=\; \frac{1}{N} \sum_i (y_i - \hat y_i)^2.
\]

Backward (per element): $2(\hat y_i - y_i)/N$ (with the
\code{'mean'} convention).

\code{nn.MSELoss(reduction='mean')}. Used for continuous-valued
regression. Implies a Gaussian likelihood with fixed variance.

\subsection{Mean Absolute Error (L1)}
\label{ssec:losses:l1}

\[
\mathcal{L}_{\text{L1}}(y, \hat y) \;=\; \frac{1}{N} \sum_i |y_i - \hat y_i|.
\]

Backward (per element): $\mathrm{sign}(\hat y_i - y_i) / N$.

\code{nn.L1Loss}. More robust to outliers than MSE because the
gradient magnitude is bounded.

\subsection{Huber loss}
\label{ssec:losses:huber}

A hybrid of L1 and L2:
\[
\mathcal{L}_{\text{Huber}}(y, \hat y; \delta) \;=\; \frac{1}{N} \sum_i
\begin{cases}
\tfrac{1}{2}(y_i - \hat y_i)^2 & |y_i - \hat y_i| \leq \delta, \\
\delta(|y_i - \hat y_i| - \tfrac{1}{2}\delta) & |y_i - \hat y_i| > \delta.
\end{cases}
\]

Quadratic near zero (so the gradient is well-behaved), linear far
from zero (so outliers do not dominate). \code{nn.HuberLoss}.

\subsection{Smooth L1}
\label{ssec:losses:smooth_l1}

A variant of Huber with $\delta = 1$ (and a slightly different
constant):
\[
\mathcal{L}_{\text{SmoothL1}} \;=\; \frac{1}{N} \sum_i
\begin{cases}
\tfrac{1}{2}(y_i - \hat y_i)^2 & |y_i - \hat y_i| < 1, \\
|y_i - \hat y_i| - \tfrac{1}{2} & \text{otherwise}.
\end{cases}
\]

\code{nn.SmoothL1Loss}. Used in object detection (Faster R-CNN's
bounding-box regression).

\section{Classification Losses}
\label{sec:losses:classification}

\subsection{Cross-entropy}
\label{ssec:losses:cross_entropy}

For a classification problem with $C$ classes, the
cross-entropy loss is
\[
\mathcal{L}_{\text{CE}}(\mathrm{logits}, y) \;=\; -\log\!\left(\frac{e^{\mathrm{logits}_y}}{\sum_c e^{\mathrm{logits}_c}}\right)
\;=\; -\mathrm{logits}_y + \mathrm{logsumexp}(\mathrm{logits}).
\]

The second form is the numerically stable computation: use
\code{logsumexp} (which has the max-shift trick built in) rather
than computing softmax then log.

\code{nn.CrossEntropyLoss(weight=None, ignore\_index=-100, label\_smoothing=0.0, reduction='mean')}.
Inputs are raw logits (not softmax outputs); the loss applies
softmax + negative-log-likelihood as a single fused op.

\subsection{The \texttt{ignore\_index} option}
\label{ssec:losses:ignore_index}

Tokens with label \code{ignore\_index} are excluded from the
loss (their per-element loss is set to 0). Used for padding
positions in language modelling: pad tokens have label
\code{-100}, and the loss skips them.

\subsection{Label smoothing}
\label{ssec:losses:label_smoothing}

A regulariser that replaces hard one-hot labels with a soft
distribution: $1 - \epsilon$ on the true class, $\epsilon /
(C-1)$ on each other class. Encourages calibrated probabilities
and discourages overconfidence:
\[
\mathcal{L}_{\text{LS}}(\mathrm{logits}, y) \;=\; -\sum_c q_c \log p_c
\]
with $q_y = 1 - \epsilon$, $q_{c \neq y} = \epsilon / (C-1)$, and
$p = \mathrm{softmax}(\mathrm{logits})$.

\subsection{Binary cross-entropy}
\label{ssec:losses:bce}

For binary classification, equivalent to softmax cross-entropy
with $C = 2$:
\[
\mathcal{L}_{\text{BCE}}(p, y) \;=\; -y \log p - (1-y) \log(1-p).
\]

\code{nn.BCELoss(p, y)} where $p$ is already in $[0, 1]$.

\code{nn.BCEWithLogitsLoss(logits, y)} accepts logits directly and
fuses the sigmoid + BCE for numerical stability:
\[
\mathcal{L} \;=\; -y \cdot \mathrm{logits} + \log(1 + e^{\mathrm{logits}})
\;=\; \max(\mathrm{logits}, 0) - y \cdot \mathrm{logits} + \log(1 + e^{-|\mathrm{logits}|}).
\]
The reformulation avoids overflow in the exponential.

\subsection{Focal loss}
\label{ssec:losses:focal}

A modification of cross-entropy \cite{lin_focal_2017} that
down-weights well-classified examples:
\[
\mathcal{L}_{\text{focal}} \;=\; -\alpha_y (1 - p_y)^\gamma \log p_y.
\]
With $\gamma = 0$ this reduces to (class-weighted) cross-entropy;
with $\gamma > 0$ the focus shifts to hard examples. Used in
object detection (RetinaNet) where the foreground/background
imbalance is extreme.

The intuition by example: with $p_y = 0.9$ (easy correct), the
modulating factor $(1 - 0.9)^2 = 0.01$ reduces the loss by
100$\times$. With $p_y = 0.1$ (hard mistake), the factor
$0.9^2 = 0.81$ barely changes the loss. The gradient signal
from hard examples therefore dominates training.

Not provided as a built-in in \code{lucid.nn} (it is a
short composite); the user can write it directly:

\begin{lstlisting}[style=lucid,language=LucidPython]
def focal_loss(logits, target, alpha=0.25, gamma=2.0):
    log_p = lucid.nn.functional.log_softmax(logits, dim=-1)
    p = log_p.exp()
    p_y = p.gather(-1, target.unsqueeze(-1)).squeeze(-1)
    log_p_y = log_p.gather(-1, target.unsqueeze(-1)).squeeze(-1)
    return -(alpha * (1 - p_y) ** gamma * log_p_y).mean()
\end{lstlisting}

\subsection{Dice and IoU losses}
\label{ssec:losses:dice_iou}

For segmentation tasks where the foreground class is sparse
(small objects in large images), cross-entropy is dominated by
the easy background pixels. The Dice loss \cite{ronneberger_unet_2015}
optimises overlap directly:
\[
  \mathcal{L}_{\text{Dice}} \;=\; 1 - \frac{2 \sum_i p_i \cdot t_i + \epsilon}{\sum_i p_i + \sum_i t_i + \epsilon},
\]
where $p_i \in [0, 1]$ is the predicted probability and
$t_i \in \{0, 1\}$ is the target. The smoothing $\epsilon$
(typically $10^{-7}$) prevents division by zero for
empty masks.

A close relative is the IoU loss
\cite{lin_focal_2017}:
\[
  \mathcal{L}_{\text{IoU}} \;=\; 1 - \frac{\sum_i p_i \cdot t_i}{\sum_i p_i + \sum_i t_i - \sum_i p_i \cdot t_i},
\]
which is differentiable and optimises a near-equivalent quantity.
Either is the typical recipe for U-Net-style biomedical
segmentation.

\section{Distribution Distances}
\label{sec:losses:distribution}

\subsection{KL divergence}
\label{ssec:losses:kl}

\code{nn.KLDivLoss(log\_target=False, reduction='batchmean')}
computes
\[
\mathcal{L}_{\text{KL}}(\mathrm{log\_p}, q) \;=\; \sum_i q_i (\log q_i - \mathrm{log\_p}_i).
\]
Input convention: the first argument is the \emph{log} of the
prediction (already log-softmaxed), the second is the target
distribution (not log).

This is the same KL discussed in
\chref{ch:distributions}~Section~\ref{sec:distributions:kl}; the
\code{nn} module wrapping is a convenience for use inside
\code{Sequential}-style training loops.

\subsection{JSD and Wasserstein}
\label{ssec:losses:jsd_wasserstein}

The Jensen-Shannon divergence and Wasserstein distances are
relevant for some GAN training procedures. \Lucid\ does not
provide these as built-in losses; users implement them when
needed.

\section{Sequence Losses}
\label{sec:losses:sequence}

\subsection{CTC loss}
\label{ssec:losses:ctc}

Connectionist Temporal Classification: a loss for alignment-free
sequence prediction where the output length differs from the
input length. Used in speech recognition (transcribe audio to
text without timestamp alignment).

\code{nn.CTCLoss(log\_probs, targets, input\_lengths, target\_lengths)}.
The mathematical content (forward-backward dynamic programming
over a special character alphabet that includes a blank symbol)
is involved; the implementation in \Lucid\ delegates to the
standard recurrence.

\subsection{Negative log-likelihood}
\label{ssec:losses:nll}

\code{nn.NLLLoss(log\_probs, targets)}: takes already-log-softmaxed
inputs and indexes into them by target. Equivalent to
\code{CrossEntropyLoss} except the user is responsible for the
log-softmax.

The split exists because some applications want to compute the
log-softmax once and use it for multiple losses or auxiliary
quantities.

\section{Pairwise and Triplet Losses}
\label{sec:losses:pairwise}

\subsection{Cosine embedding}
\label{ssec:losses:cosine}

\code{nn.CosineEmbeddingLoss(x1, x2, y, margin=0)}: $y = 1$
pushes $x_1$ and $x_2$ to be cosine-similar, $y = -1$ pushes
them to be dissimilar:
\[
\mathcal{L} \;=\; \begin{cases}
1 - \cos(x_1, x_2) & y = 1, \\
\max(0, \cos(x_1, x_2) - m) & y = -1.
\end{cases}
\]

\subsection{Triplet margin}
\label{ssec:losses:triplet}

\code{nn.TripletMarginLoss(anchor, positive, negative, margin=1.0)}:
push the anchor closer to the positive than to the negative by at
least the margin:
\[
\mathcal{L} \;=\; \max(0, \|a - p\|_2 - \|a - n\|_2 + m).
\]
Used in metric learning (face recognition, image retrieval).

\section{Numerical Stability}
\label{sec:losses:numerical_stability}

The most common loss-related bug is the unstable composition of
softmax then negative log:
\begin{lstlisting}[style=lucid,language=LucidPython]
# UNSTABLE: softmax can underflow to 0, log(0) = -inf
p = logits.softmax(dim=-1)
loss = -(p * target).log().sum()
\end{lstlisting}

The stable form fuses the two:
\begin{lstlisting}[style=lucid,language=LucidPython]
# STABLE: log-softmax via logsumexp
log_p = logits.log_softmax(dim=-1)
loss = -(log_p * target).sum()
\end{lstlisting}

\code{CrossEntropyLoss} and \code{BCEWithLogitsLoss} encapsulate
the stable fusion. The rule of thumb: if the input is logits,
use the \code{*WithLogitsLoss} variant.

\subsection{The log-sum-exp trick}
\label{ssec:losses:logsumexp}

The log-sum-exp identity, the foundation of stable softmax:
\[
  \log \sum_i e^{x_i} \;=\; \max_i x_i + \log \sum_i e^{x_i - \max_i x_i}.
\]
Subtracting the max from every $x_i$ shifts the largest
exponent to zero (so $e^0 = 1$ is the largest term and no
overflow occurs); the sum is then between 1 and $n$ (where $n$
is the number of terms), and its log is bounded by $\log n$.
The subtracted max is added back outside the log, preserving
the value.

The cross-entropy loss using log-sum-exp:
\[
  -\log \frac{e^{x_y}}{\sum_c e^{x_c}}
  \;=\;
  \log \sum_c e^{x_c} - x_y
  \;=\;
  \max_c x_c + \log \sum_c e^{x_c - \max_c x_c} - x_y.
\]

\subsection{Overflow vs underflow regimes}
\label{ssec:losses:overflow_underflow}

For FP32, the safe range of $e^x$ is approximately
$x \in [-87, 88]$ (outside which we get underflow or overflow).
For FP16, the range narrows to roughly $[-15, 11]$ --- which is
why softmax and the log-loss combinations are blacklisted
under FP16 AMP (\chref{ch:amp}); they are silently promoted
to FP32 for the unstable steps and cast back for the result.

The log-sum-exp trick is always applied; the dtype promotion is
an additional safety belt for the tail of cases where even the
shifted log-sum-exp overflows. In practice, both are needed:
the trick eliminates the common case, the FP32 promotion
catches the rare ones.

\section{Contrastive Losses}
\label{sec:losses:contrastive}

Self-supervised learning has popularised a family of losses
based on contrasting pairs of representations.

\subsection{NT-Xent (SimCLR)}
\label{ssec:losses:nt_xent}

For a batch of $2N$ representations consisting of $N$ positive
pairs $(z_i, z_{i^+})$, the NT-Xent loss is
\[
  \mathcal{L} \;=\;
  -\log \frac{\exp(\mathrm{sim}(z_i, z_{i^+}) / \tau)}
              {\sum_{k \neq i} \exp(\mathrm{sim}(z_i, z_k) / \tau)},
\]
where $\mathrm{sim}$ is cosine similarity and $\tau$ is the
temperature (typically 0.07--0.5). The positive pair is two
augmented views of the same image; the negatives are all other
samples in the batch.

The temperature controls the loss landscape's sharpness: small
$\tau$ heavily penalises near-positive negatives; large $\tau$
treats them more uniformly. Typical default $\tau = 0.1$.

\subsection{InfoNCE}
\label{ssec:losses:infonce}

A more general formulation:
\[
  \mathcal{L}_{\text{InfoNCE}}
  \;=\;
  -\log \frac{\exp(f(x, c) / \tau)}{\sum_{x' \in \mathcal{N} \cup \{x\}} \exp(f(x', c) / \tau)},
\]
where $f(x, c)$ is a similarity score between a query $x$ and
a context $c$, and $\mathcal{N}$ is a set of negative samples.
NT-Xent is a special case with $c = z_{i^+}$ and the in-batch
negatives.

InfoNCE is the loss behind CLIP (text--image pairs), MoCo
(momentum-encoded negatives), and many other contrastive
methods.

\section{Class Imbalance}
\label{sec:losses:class_imbalance}

For datasets where one class dominates (medical diagnosis,
fraud detection), the standard cross-entropy gives the majority
class disproportionate gradient. Three mitigations:

\subsection{Per-class weight}
\label{ssec:losses:class_weight}

\code{CrossEntropyLoss(weight=w)} where $w \in \mathbb{R}^C$ scales
each class's loss contribution. Setting $w_c = 1 / N_c$
(inverse frequency) equalises the average per-class gradient.

\subsection{Focal loss}
\label{ssec:losses:focal_imbalance}

Down-weights easy examples (typically majority-class
correct-prediction examples). Discussed above.

\subsection{Re-sampling}
\label{ssec:losses:resampling}

Re-sample the dataset to balance classes (oversample minority,
undersample majority). Handled at the data-loader level, not the
loss level.

\section{Multi-Task Losses}
\label{sec:losses:multi_task}

A multi-task model has multiple outputs and multiple losses.
The standard pattern is to sum the losses with task-specific
weights:
\[
\mathcal{L} \;=\; \sum_t \lambda_t \mathcal{L}_t.
\]

The weights $\lambda_t$ balance the task contributions. They are
usually hand-tuned; learnable variants (uncertainty weighting,
GradNorm) exist but are not built into \Lucid.

\section{Performance and Dispatch}
\label{sec:losses:perf}

Most loss functions are sequences of primitives (matmul, softmax,
reduction). They dispatch through the standard backend with no
special fast paths.

\code{CrossEntropyLoss} is the exception: the fused
softmax-cross-entropy-with-logits kernel is in the 3.4 per-op
\MPSGraph\ carve-out
(\chref{ch:mpsgraph_dispatch}). The fused kernel computes the
loss and the gradient in one pass, saving the intermediate
softmax materialisation.

\tabref{tab:losses:perf} reports measured times on M3~Max.

\begin{table}[t]
\centering
\small
\begin{tabular}{lrr}
\toprule
Loss & Tensor shape & Time (ms) \\
\midrule
MSE                       & $(B, D) = (256, 1024)$ & 0.05 \\
L1                        & $(256, 1024)$          & 0.06 \\
Huber                     & $(256, 1024)$          & 0.08 \\
Cross-entropy (fused)     & logits $(B, C) = (256, 50000)$ & 0.45 \\
Cross-entropy (composite) & same                     & 2.1  \\
BCEWithLogitsLoss         & $(B, D) = (256, 1)$    & 0.04 \\
KLDivLoss                 & $(256, 1024)$          & 0.18 \\
CTCLoss                   & seq=$(T, B, C) = (200, 32, 30)$ & 8.2 \\
\bottomrule
\end{tabular}
\caption[Loss function timings.]{Per-call timings for
representative loss functions on M3~Max. The cross-entropy fused
path is 4--5$\times$ faster than the equivalent composite. CTC's
high cost is the dynamic-programming forward/backward over the
sequence.}
\label{tab:losses:perf}
\end{table}

\section{Choosing a Loss}
\label{sec:losses:choosing}

A short decision guide for new tasks.

\begin{itemize}
  \item \textbf{Regression of continuous values}: MSE (default),
    L1 (robust), Huber (compromise).
  \item \textbf{Multi-class classification}: CrossEntropyLoss with
    raw logits.
  \item \textbf{Multi-label classification (multiple binary
    decisions per sample)}: BCEWithLogitsLoss.
  \item \textbf{Variational autoencoder}: MSE / BCE for
    reconstruction + KL divergence to prior. See
    \chref{ch:distributions}.
  \item \textbf{Speech recognition or other alignment-free
    sequence prediction}: CTCLoss.
  \item \textbf{Metric learning (face recognition, retrieval)}:
    TripletMarginLoss or CosineEmbeddingLoss.
  \item \textbf{Object detection}: SmoothL1Loss for bounding
    boxes + CrossEntropyLoss (or focal) for class.
\end{itemize}

\section{Implementation Notes}
\label{sec:losses:impl_notes}

Loss modules live in \code{lucid/nn/modules/loss.py}; the
functional forms in \code{lucid/nn/functional/loss.py}.
Most are short composites; the only C++-primitive losses are
CrossEntropy (fused with softmax) and CTC (the dynamic-programming
recurrence is impractical in pure Python at production speeds).

\section{Summary and Forward References}
\label{sec:losses:summary}

\Lucid\ provides the standard loss family at \code{lucid.nn} and
\code{lucid.nn.functional}, covering regression (MSE / L1 / Huber
/ SmoothL1), classification (CrossEntropy / BCE / focal),
distribution distances (KL), and sequence losses (CTC, NLL). The
\code{*WithLogitsLoss} variants and the fused softmax-cross-
entropy are the numerically stable choices when the input is raw
logits.

The next chapter, \chref{ch:initialization}, treats weight
initialisation theory. \chref{ch:optimizers} treats optimisers.
\chref{ch:amp} treats mixed precision.

\begin{lucidnote}
For the user: the loss API is drop-in compatible with
\refframework. The most common choice is
\code{nn.CrossEntropyLoss()} for classification and
\code{nn.MSELoss()} for regression; both expect raw outputs
(logits or values) and handle the appropriate fusion internally.
\end{lucidnote}
